\chapter{First Aid Practice Scenarios}\label{ch:scenarios}

Apply what you have learned. Each scenario describes a realistic situation. Think through your response before reading the suggested action, then compare. Cover the answer with a card until you have decided.

\section{Scenario 1: Unresponsive Adult at Home}

You enter the kitchen and find your 68-year-old neighbor lying on the floor. He does not respond when you call his name or tap his shoulder. His chest is moving with irregular gasping sounds.

\textbf{What do you do?}

\textbf{Response:} Check scene safety first. Call emergency services immediately (or have a bystander call) and retrieve an AED if one is available. The gasping sounds are agonal breathing, not normal breathing. Begin chest compressions immediately: hands on the center of the chest, at least 2 inches (5 cm) deep, 100--120 per minute. If you are trained and willing, give 2 rescue breaths after every 30 compressions. Continue until an AED arrives, EMS takes over, or the person shows signs of life.

\section{Scenario 2: Choking at a Restaurant}

At dinner, a man at the next table suddenly stands up, clutches his throat with both hands, and makes no sound. His lips are turning blue.

\textbf{What do you do?}

\textbf{Response:} This is complete airway obstruction -- the universal choking sign (hands to throat) with no sound. Ask loudly "Are you choking?" He cannot answer, so act immediately: perform abdominal thrusts (Heimlich maneuver) -- stand behind him, wrap your arms around his waist, make a fist just above the navel, grasp it with the other hand, and thrust sharply inward and upward until the object comes out or he becomes unresponsive. If he becomes unresponsive, lower him to the floor, call emergency services, and begin CPR, checking the mouth for the object before giving breaths.

\section{Scenario 3: Severe Leg Bleeding}

A kitchen accident has cut the front of your friend's thigh. Blood is bright red and spurting with each heartbeat, pooling rapidly on the floor.

\textbf{What do you do?}

\textbf{Response:} This is arterial bleeding -- life-threatening. Call emergency services immediately. Put on gloves if available. Press directly on the wound with a clean cloth or dressing with firm, continuous pressure. Do not remove the dressing when it soaks through -- add more layers on top. If bleeding continues despite firm pressure, apply a tourniquet 2--3 inches above the wound (above the bleeding site, not over a joint), tighten until the bleeding stops, note the time clearly, and do not loosen it. Keep the person lying down, warm, and monitored for shock.

\section{Scenario 4: Boiling Water Burn}

A child pulls a pot of boiling water onto herself. Her forearm is red, blistered, and very painful.

\textbf{What do you do?}

\textbf{Response:} This is a partial-thickness (second-degree) burn. Run cool -- not cold -- water over the burn for 10--20 minutes as soon as possible. Do not use ice. Do not break blisters, peel skin, or apply butter, oils, or ointments. Cover loosely with a clean, dry, non-stick dressing or clean cloth. Since a child's smaller burn can be significant, seek medical evaluation. Monitor for shock and keep her calm and warm.

\section{Scenario 5: Chest Pain at the Mall}

A 60-year-old man at the mall leans against a wall, clutching his chest, pale and sweating. He says the pressure has been building for 15 minutes and he feels short of breath.

\textbf{What do you do?}

\textbf{Response:} Suspect a heart attack. Have him sit down, rest, and loosen tight clothing. Call emergency services immediately. Ask if he has aspirin and whether he is allergic or has been advised by his doctor against it; if safe, have him chew one adult (325 mg) or four low-dose (81 mg) aspirins. Ask about a history of heart disease and whether he takes nitroglycerin (help him use his own if prescribed). Stay calm, reassure him, and monitor responsiveness and breathing. If he becomes unresponsive and stops breathing normally, begin CPR and use an AED.

\section{Scenario 6: Possible Stroke}

Your grandfather, who speaks normally, suddenly has slurred speech and his smile is crooked. He cannot raise his right arm. He says it started "about an hour ago."

\textbf{What do you do?}

\textbf{Response:} Use FAST: Face drooping, Arm weakness, Speech difficulty -- Time to call emergency services. Call immediately; note the exact time symptoms began or, if unknown, the "last known well" time -- it determines eligibility for time-critical treatment. Do not give him food, drink, or aspirin (a stroke may be bleeding). Keep him calm, lying on his side if he is drowsy, and do not let him walk. Note any change in symptoms to report to responders.

\section{Scenario 7: Shaky, Confused Friend}

Your friend with type 1 diabetes becomes shaky, sweaty, confused, and irritable. He forgot to eat lunch.

\textbf{What do you do?}

\textbf{Response:} This is likely hypoglycemia (low blood sugar). He is conscious, so give fast-acting sugar immediately: glucose tablets, juice, regular soda, or sugar dissolved in water. If he improves in 10--15 minutes, follow with a snack containing longer-acting carbohydrate. If he does not improve, call emergency services. If he becomes unconscious, do NOT give food or drink (choking risk) -- call emergency services, place him in the recovery position, and monitor breathing.

\section{Scenario 8: Seizure in a Shopping Aisle}

A person collapses in a shopping aisle, stiffens, and begins convulsing. Shoppers gather around; someone suggests putting a spoon in the mouth.

\textbf{What do you do?}

\textbf{Response:} Do NOT restrain the person and never put anything in their mouth -- it can cause choking or broken teeth. Protect them: clear hard objects, cushion the head, and remove eyeglasses. Time the seizure. When the convulsions stop, turn them onto their side (recovery position), check breathing, and stay with them. Call emergency services if the seizure lasts longer than 5 minutes, a second seizure follows without recovery, breathing is difficult afterward, they are injured, or it is their first seizure. Politely stop bystanders from interfering.

\section{Scenario 9: Bee Sting with Throat Swelling}

A friend is stung by a bee and within minutes develops hives, swelling of the face, and a hoarse voice, and says it is hard to breathe.

\textbf{What do you do?}

\textbf{Response:} This is anaphylaxis -- a life-threatening emergency. Call emergency services immediately. If the person has an epinephrine auto-injector, help them use it: inject into the outer thigh muscle, through clothing if necessary, and hold for several seconds. Keep them lying down (sitting up if breathing is easier) and elevate the legs. Monitor responsiveness and breathing closely; be ready to begin CPR if they become unresponsive and stop breathing. A second epinephrine dose may be given after 5--15 minutes if symptoms persist and you are trained. Even if they improve, EMS must still evaluate them -- reactions can recur.

\section{Scenario 10: Child Pulled from a Pool}

A 5-year-old is pulled unresponsive from a backyard pool. She is not breathing.

\textbf{What do you do?}

\textbf{Response:} This is a drowning -- call emergency services immediately (or direct a bystander). Begin CPR right away: give 5 initial rescue breaths, then cycles of 30 compressions to 2 breaths (compressions about 2 inches/5 cm deep, 100--120 per minute), allowing water to drain with gravity during positioning. Use an AED as soon as available. Continue until the child shows signs of life or EMS takes over. Children can recover remarkably well from drowning, so persist with CPR even if it seems long.

\section{Scenario 11: Hiker Shivering Uncontrollably}

A hiker on a cold, windy ridge is shivering hard, clumsy, and speaking slowly. His clothes are wet.

\textbf{What do you do?}

\textbf{Response:} This is moderate hypothermia. Get him out of the wind and cold immediately: into a tent, behind a windbreak, or inside a shelter. Replace wet clothes with dry layers and add insulation. Give warm (not hot) sweet drinks if he is fully awake and able to swallow. Do NOT rub his limbs, give alcohol, or warm his extremities directly (risk of "afterdrop" arrhythmia). Handle him gently. If he becomes confused, stops shivering, or loses consciousness, treat as severe hypothermia: handle with extreme care, protect from further heat loss, and arrange urgent evacuation -- do not attempt rapid field rewarming you cannot sustain.

\section{Scenario 12: Opioid Overdose in a Restroom}

A person is found slumped in a public restroom, unresponsive, with very slow, shallow breathing and pinpoint pupils.

\textbf{What do you do?}

\textbf{Response:} Call emergency services immediately -- this is likely opioid overdose and it is reversible. If naloxone (Narcan) is available and you are trained, give one dose into the nose or muscle and repeat every 2--3 minutes if there is no response. Give rescue breathing (1 breath every 5--6 seconds for adults) if breathing is absent or inadequate. Place the person in the recovery position if breathing. Stay with them: naloxone wears off in 30--90 minutes and a second overdose can occur. Even if they wake up, they must be evaluated by medical personnel.

\section{Scenario 13: Fallen from a Ladder}

A man has fallen from a ladder and lies still on the ground. He is conscious but cannot feel his legs. Neighbors want to carry him inside.

\textbf{What do you do?}

\textbf{Response:} Suspect spinal injury. Do NOT let anyone move or lift him. Call emergency services. Ask one person to hold his head and neck still (manual in-line stabilization) while you assess breathing and responsiveness. Do not move him unless there is immediate danger (fire, traffic, collapsing structure) -- and if you must move him, do it with several helpers rolling him as a unit with the head and neck held in line. Keep him still and warm, and monitor breathing until professional rescuers arrive with proper immobilization.

\section{Scenario 14: Jellyfish Sting at the Beach}

A swimmer runs out of the water in pain, with red, stinging welts across her calf from a jellyfish.

\textbf{What do you do?}

\textbf{Response:} Rinse the area with seawater -- do NOT use fresh water (it can trigger more stinging). Remove any visible tentacles with tweezers or a gloved hand, not bare skin. Do not rub the area or rinse with vinegar (vinegar is reserved for box jellyfish stings). Soak the area in hot water (up to 45\textdegree C / 113\textdegree F) or apply a cold pack to relieve pain, then treat the pain with oral analgesics. Watch for signs of a severe reaction -- breathing difficulty, chest pain, dizziness, or swelling elsewhere -- and call emergency services if any develop.

\section{Scenario 15: Active Bleeding from a Dog Bite}

A neighbor's dog bites a delivery worker on the forearm, leaving a deep puncture wound that is bleeding steadily.

\textbf{What do you do?}

\textbf{Response:} Put on gloves. Control bleeding with firm direct pressure. Wash the wound gently with clean water and mild soap (puncture wounds are deep and easily infected). Cover with a clean dressing. Advise the person to seek medical care: animal bites carry high infection risk, may require antibiotics, and tetanus status must be checked. Note the animal's owner and description -- if the animal is stray, wild, or acting strangely, report to animal control for rabies evaluation.

\section{Scenario 16: Collapsed Runner in Midday Heat}

At a summer race, a runner collapses near the finish line. Her skin is hot and dry, she is confused, and her pulse is racing. She is not sweating.

\textbf{What do you do?}

\textbf{Response:} This is heat stroke -- a life-threatening emergency. Call emergency services immediately. Begin rapid cooling at once: move her to shade, remove excess clothing, and immerse her in cold water if available, or pour/spray cool water over her and fan aggressively; place ice packs on her neck, armpits, and groin. Do NOT give fluids if she cannot swallow safely. Monitor her breathing and consciousness continuously and be ready to begin CPR if she becomes unresponsive. Even if she seems to recover, she must be evaluated at a hospital.

\section{Scenario 17: Chemical Splash in the Eye}

A splash of drain cleaner lands in a gardener's eye. He is in severe pain and blinking hard.

\textbf{What do you do?}

\textbf{Response:} This is an alkali chemical eye injury -- an emergency. Flush the eye immediately with clean water or saline for at least 20--30 minutes (alkali continues to damage tissue). Hold his eyelids open, tilt his head so the water runs from the nose outward, and keep flushing continuously even while preparing transport. Do not let him rub the eye or neutralize the chemical with anything else. Call emergency services and transport to hospital care during or immediately after flushing.

\section{Scenario 18: Glass Embedded in the Forearm}

A coworker falls onto a broken jar, and a large piece of glass is stuck in her forearm, with blood soaking her sleeve around the wound.

\textbf{What do you do?}

\textbf{Response:} Do NOT remove the glass -- it may be plugging the wound and controlling bleeding, and removing it can cause severe new bleeding. Call emergency services. Control bleeding around the object with firm pressure on the surrounding skin using dressings, then build bulky dressings around the glass to stabilize it in place so it cannot move. Do not press directly on the glass itself. Keep her still and monitor for shock until paramedics arrive.

\section{Scenario 19: Unconscious Person with Diabetes}

You find a man unresponsive on the floor at the gym, breathing slowly. A friend says he has type 1 diabetes and may have skipped a meal before his workout.

\textbf{What do you do?}

\textbf{Response:} Suspect a diabetic emergency (likely severe hypoglycemia), but you cannot confirm it -- call emergency services immediately. Do NOT give food or drink to an unconscious person (choking risk). Place him in the recovery position to keep the airway clear and protect him if he vomits. Monitor breathing and responsiveness. If you are trained and it is available, follow the person's diabetes emergency plan or glucose gel instructions (glucose gel can be applied to the inside of the cheek of an unconscious person only if the care plan directs it and he can swallow safely -- otherwise, do not). Report everything to responders.

\section{Scenario 20: Distressed Swimmer the Day After}

A child who was rescued from a pool yesterday and seemed fine is now drowsy, coughing, and struggling to breathe.

\textbf{What do you do?}

\textbf{Response:} Suspect secondary (dry) drowning -- water irritation in the lungs causing symptoms hours later. Call emergency services immediately. Keep the child sitting upright if that eases breathing. Do not give food or drink. Reassure the child and monitor breathing continuously until help arrives. Even delayed respiratory symptoms after any water submersion need urgent medical evaluation.

\section{Review Questions}

\begin{enumerate}[leftmargin=*]
\item In Scenario 1, why are "gasping sounds" treated as not breathing?
\item Why must a blood-soaked dressing never be removed while controlling bleeding?
\item In Scenario 6, why must you record the exact time symptoms began?
\item In Scenario 8, why must nothing ever be placed in a seizing person's mouth?
\item In Scenario 13, why must the fallen person not be moved or lifted?
\item In Scenario 12, why must the person remain with medical personnel even after naloxone wakes them?
\item In Scenario 16, why is rapid cooling the top priority even before transport begins?
\item In Scenario 18, why must the embedded glass never be removed?
\item In Scenario 20, why does a child who seemed fine after a rescue still need urgent care?
\end{enumerate}
\index{practice scenarios}
\index{case studies}
\index{scenario-based learning}
